\documentclass[11pt]{article}
\usepackage[T1]{fontenc}
\usepackage[polish]{babel}
\usepackage{amsmath}
%\usepackage{amssymb}
%\topmargin -2cm
%\textheight 23cm
%\oddsidemargin -1cm
%\evensidemargin -1cm
%\textwidth 17.5cm
\usepackage[lmargin=1cm,rmargin=1.5cm,tmargin=1cm,bmargin=1cm]{geometry}
\usepackage{natbib}
\usepackage{graphicx}

\title{Zadania domowe na środę 11 marca 2020}
\author{Mateusz Trubiłowicz}

\begin{document}

\maketitle
\medskip

\begin{enumerate}

%\begin{center}
%Zadania domowe na środę 11 marca 2020r. \\
%Mateusz Trubiłowicz
%\end{center}
%POWYŻSZE ODKOMENTOWAĆ, JEŚLI CHCE SIĘ MNIEJ KRZYKLIWY TYTUŁ, BEZ DATY

\item
\begin{Large}
Zadanie 1.
\end{Large}
\par
Tutaj wpisuję treść zadania, które mam rozwiązać. Jeśli w środku pojawiają się wyrażenia symboliczne typu \( G \times H\), używam  \(\backslash( \) i \(\backslash) \) %Czyli \( i \) 
aby wejść w "tryb matematyczny".\par

\medskip
\textbf{Rozwiązanie: }

W treści rozwiązania, jeśli chcę, aby pojawiło się równanie w środku tekstu, wykonuję podobną czynność. Jesli chcę, by równanie stanowiło osobną linijkę robię:

\[ [g^{-1},h^{-1}]=ghg^{-1}h^{-1}=h'\in H \] 

(nawiasy okrągłe zastępuję kwadratowymi).


\item
\begin{Large}
Zadanie 2.
\end{Large}
\par

Tutaj ponownie wpisuję główną treść zadania\par

a) Tutaj z kolei opisuję poszczególne podpunkty. Bardzo często używam \(trybu \ matematycznego\).\par
\medskip
\textbf{Rozwiązanie: } 
Tutaj odpowiadam na pierwszy podpunkt. Jeśli w rozwiązaniu pojawia się jakies równanie, do którego chcę się odwołać, typu:

\begin{equation}
\label{PODPIS}  %%TUTAJ JEST PODPIS
\begin{cases}
\hfill \frac{d \epsilon(\mathbf{x},t)}{d t} =& -k(u) \epsilon \\
\hfill \frac{\partial u(\mathbf{x},t)}{\partial t} - 
u \Delta u(\mathbf{x},t) =& f(\mathbf{x},t) 
\end{cases} \quad\quad  (\mathbf{x},t) \in \Omega \times (0,T],
\end{equation}

Używam begin{equation} i labeluję. Odwołać do równania (\ref{PODPIS}) mogę się używając komendy jak obok.

\par
\bigskip

b) Kolejna treść...\par
\medskip
\textbf{Rozwiązanie: }
Przykład macierzy (\ref{macierz}), do której potem będę się chciał odwołać:
\begin{equation}
\label{macierz}
    A = \left(
    \begin{array}{ccc} %liczba liter "c" odpowiada liczbie kolumn
    3 & -1 & 2 \\ 	
    0 & 1 & 2 \\ 
    1 & 0 & -1 
\end{array} 
\right).  
\end{equation}

\end{enumerate}

\end{document}
